\section{Background}
\begin{frame}{Background: On-Policy vs.\ Off-Policy}
\begin{itemize}
    \item \textbf{On-policy learning:} use the deterministic outcomes or samples from the target policy to train the algorithm
    \begin{itemize}
        \item has low sample efficiency (TRPO, PPO, A3C)
        \item require new samples to be collected for nearly every update to the policy
        \item becomes extremely expensive when the task is complex
    \end{itemize}
    \item \textbf{Off-policy methods:} training on a distribution of transitions or episodes produced by a different behavior policy rather than that produced by the target policy
    \begin{itemize}
        \item Does not require full trajectories and can reuse any past episodes (experience replay) for much better sample efficiency
        \item relatively straightforward for Q-learning based methods
    \end{itemize}
\end{itemize}
\end{frame}
\begin{frame}{Background: Bellman Equation}
\begin{itemize}
    \item \textbf{Value Function:} How good is a state?
\end{itemize}
\begin{align*}
V(s) &= \mathbb{E}[G_t \mid S_t = s] \\
     &= \mathbb{E}[R_{t+1} + \gamma R_{t+2} + \gamma^2 R_{t+3} + \ldots \mid S_t = s] \\
     &= \mathbb{E}[R_{t+1} + \gamma (R_{t+2} + \gamma R_{t+3} + \ldots) \mid S_t = s] \\
     &= \mathbb{E}[R_{t+1} + \gamma G_{t+1} \mid S_t = s] \\
     &= \mathbb{E}[\underbrace{R_{t+1} + \gamma V(S_{t+1})}_{\textcolor{red}{\text{temporal difference target}}} \mid S_t = s]
\end{align*}
\begin{itemize}
    \item Similarly, for \textbf{Q-Function:} How good is a state-action pair?
\end{itemize}
\begin{align*}
Q(s,a) &= \mathbb{E}[R_{t+1} + \gamma V(S_{t+1}) \mid S_t = s, A_t = a] \\
       &= \mathbb{E}[R_{t+1} + \gamma\, \mathbb{E}_{a \sim \pi} Q(S_{t+1}, a) \mid S_t = s, A_t = a]
\end{align*}
\end{frame}
\begin{frame}{Background: Value-Based Method}
\begin{itemize}
    \item $\ldots, S_t, A_t, R_{t+1}, S_{t+1}, A_{t+1}, \ldots$ \textbf{(on-policy):}
\end{itemize}
\[
Q(S_t, A_t) \leftarrow Q(S_t, A_t) + \alpha\big( R_{t+1} + \gamma Q(S_{t+1}, A_{t+1}) - Q(S_t, A_t) \big)
\]
\begin{itemize}
    \item \textbf{Q-Learning (off-policy)}
\end{itemize}
\[
Q(S_t, A_t) \leftarrow Q(S_t, A_t) + \alpha\big( R_{t+1} + \gamma \max_{a \in \mathcal{A}} Q(S_{t+1}, a) - Q(S_t, A_t) \big)
\]
\begin{itemize}
    \item \textbf{DQN}, Mnih et al., 2015
\end{itemize}
\[
\mathcal{L}(\theta) = \mathbb{E}_{(s,a,r,s') \sim U(D)}\!\left[ \Big( r + \gamma \max_{a'} Q(s', a'; \theta^{-}) - Q(s, a; \theta) \Big)^{2} \right]
\]
\begin{itemize}
    \item Function Approximation
    \item Experience Replay: samples randomly drawn from replay memory
    \item Doesn't scale to continuous action space
\end{itemize}
\end{frame}
\begin{frame}{Background: Policy-Based Method (Actor-Critic)}
\begin{itemize}
    \item \textbf{Critic:} updates value function parameters $w$ and depending on the algorithm it could be action-value $Q(a \mid s; w)$ or state-value $V(s; w)$.
    \item \textbf{Actor:} updates policy parameters $\theta$, in the direction suggested by the critic, $\pi(a \mid s; \theta)$.
\end{itemize}
\begin{align*}
\nabla \mathcal{J}(\theta) &= \mathbb{E}_{\pi_\theta}\!\left[ \nabla \ln \pi(a \mid s, \theta)\, Q_\pi(s, a) \right] && \text{policy gradient} \\
\theta &\leftarrow \theta + \alpha_\theta\, Q(s, a; w)\, \nabla_\theta \ln \pi(a \mid s; \theta) && \text{update actor} \\
G_{t:t+1} &= r_t + \gamma Q(s', a'; w) - Q(s, a; w) && \text{correction for action-value} \\
w &\leftarrow w + \alpha_w\, G_{t:t+1}\, \nabla_{\!w} Q(s, a; w) && \text{update critic}
\end{align*}
\end{frame}
